\section{Lambda Expressions}

\subsection{Section Overview}
Lambda expressions, introduced in C++11, provide a convenient way to define anonymous function objects right at the location where they are needed. They are particularly useful when working with STL algorithms that require a function as an argument. Lambdas make code more concise and readable by keeping the operation logic close to where it's used.

\subsection{Motivation for Lambdas}
Before lambdas, if you needed a simple, one-off function for an STL algorithm (like \texttt{std::for\_each}), you had to either define a full function elsewhere or create a function object (a class that overloads \texttt{operator()}). This could be verbose and separate the logic from its use. Lambdas allow you to define this logic ``inline''.

\subsection{Structure of a Lambda Expression}
A lambda expression has the following structure:
\begin{center}
\texttt{[capture\_list](parameter\_list) -> return\_type \{ function\_body \}}
\end{center}
\begin{itemize}
    \item \textbf{\texttt{[capture\_list]}}: Specifies which variables from the enclosing scope are accessible inside the lambda and how they are captured (by value or by reference).
    \item \textbf{\texttt{(parameter\_list)}}: The parameters that the lambda takes, just like a regular function.
    \item \textbf{\texttt{-> return\_type}}: The return type of the lambda. This is often optional, as the compiler can deduce it.
    \item \textbf{\texttt{\{ function\_body \}}}: The code that the lambda executes.
\end{itemize}

\begin{lstlisting}
// A simple lambda that takes two integers and returns their sum
auto my_lambda = [](int a, int b) -> int { return a + b; };
std::cout << my_lambda(5, 3); // 8
\end{lstlisting}

\subsection{Coding Exercise 43: Using Lambda Expressions --- Exercise 1}
Write a lambda expression that takes two doubles and returns their product.
\begin{lstlisting}
#include <iostream>

int main() {
	auto product = [](double a, double b) { return a * b; };
	std::cout << "Product of 3.5 and 4.0 is " << product(3.5, 4.0) << std::endl;
	return 0;
}
\end{lstlisting}

\subsection{Stateful Lambdas}
The capture list allows a lambda to ``capture'' variables from its surrounding scope, making it ``stateful''.
\begin{itemize}
    \item \texttt{[]}: Capture nothing.
    \item \texttt{[x, \&y]}: Capture \texttt{x} by value and \texttt{y} by reference.
    \item \texttt{[=]}: Capture all external variables by value.
    \item \texttt{[\&]}: Capture all external variables by reference.
    \item \texttt{[=, \&y]}: Capture \texttt{y} by reference and everything else by value.
\end{itemize}
\begin{lstlisting}
int x = 10;
int y = 20;

auto my_lambda = [x, &y]() {
	// x is a copy, y is the original
	std::cout << "x: " << x << ", y: " << y << std::endl;
	y = 30; // this changes the original y
};
my_lambda();
std::cout << "Original y is now: " << y << std::endl; // 30
\end{lstlisting}

\subsection{The Standard Library Functional Header}
The \texttt{<functional>} header provides \texttt{std::function}, a general-purpose polymorphic function wrapper. \texttt{std::function} objects can store, copy, and invoke any callable target---functions, lambda expressions, bind expressions, or other function objects.
\begin{lstlisting}
#include <functional>

void perform_operation(int a, int b, std::function<int(int, int)> op) {
	std::cout << "Result: " << op(a, b) << std::endl;
}

int main() {
	auto add = [](int a, int b){ return a + b; };
	perform_operation(10, 5, add); // Result: 15
}
\end{lstlisting}

\subsection{Section Challenge}
Use a lambda expression with the \texttt{std::for\_each} algorithm to square each element in a vector of integers.
\subsection{Section Challenge --- Solution}
\begin{lstlisting}
#include <iostream>
#include <vector>
#include <algorithm>

int main() {
	std::vector<int> vec {1, 2, 3, 4, 5};

	std::for_each(vec.begin(), vec.end(), [](int &x) {
		x = x * x; // square the element
	});

	for (int x : vec) {
		std::cout << x << " "; // 1 4 9 16 25
	}
	std::cout << std::endl;

	return 0;
}
\end{lstlisting}

\subsection{Quiz 18: Section 21 Quiz}
\begin{enumerate}
    \item What is a primary benefit of using lambda expressions?
    \begin{enumerate}
        \item[a)] They make the program run significantly faster.
        \item[b)] They allow you to define anonymous function objects inline, right where they are needed.
        \item[c)] They are the only way to work with STL algorithms.
        \item[d)] They replace the need for classes.
    \end{enumerate}

    \item In a lambda \texttt{[=, \&x]()\{\}}, what does the capture list mean?
    \begin{enumerate}
        \item[a)] Capture all variables by reference, except \texttt{x} which is by value.
        \item[b)] Capture all variables by value, except \texttt{x} which is by reference.
        \item[c)] Capture only \texttt{x} by value.
        \item[d)] This is a syntax error.
    \end{enumerate}

    \item The \texttt{auto} keyword is often used with lambdas. Why?
    \begin{enumerate}
        \item[a)] It is required by the lambda syntax.
        \item[b)] The type of a lambda is unique and unnamed, so \texttt{auto} is the most convenient way to store it in a variable.
        \item[c)] It makes the lambda stateful.
        \item[d)] It specifies the return type of the lambda.
    \end{enumerate}
\end{enumerate}

\textbf{Answers:} 1. (b), 2. (b), 3. (b)
